% !PRINTED 205
\weberpage{205}

%<<CONT
\[
\frac{1-x^2}{1-\varkappa^2 x^2}, \qquad \frac{1-\varkappa^2 x^2}{\varkappa^2(1-x^2)}, \qquad \frac{1}{\varkappa^2 x^2}
\]

the roots of $B = 0$, $C = 0$, $D = 0$. The meanings of these latter
roots, however, are % [source typo] German prints the singular although three roots precede.

\[
\operatorname{sn}\left(\frac{(4\mu+1)K + 4\mu' iK'}{n}\right)^2,
\]
\[
\operatorname{sn}\left(\frac{(4\mu+1)K + (4\mu'+1)iK'}{n}\right)^2,
\]
\[
\operatorname{sn}\left(\frac{4\mu K + (4\mu'+1)iK'}{n}\right)^2,
\]

and accordingly, by the solution of the equation $A = 0$, all
quantities of the form

\[
\operatorname{sn}\left(\frac{\mu K + \mu' iK'}{n}\right)^2,
\]

in which $\mu$, $\mu'$ are arbitrary integers, are \gesperrt{rationally determined}.

\begin{center}
\S\ 62. \quad \gesperrt{The Abelian Relations.}
\end{center}

For a closer investigation of the algebraic nature of the equation of
the division of the periods, a system of relations among its roots is
of great importance, and we now pass on to the derivation of that
system \textsuperscript{1}).

We consider the sum:

\weq{1}{\sum^{\nu} e^{\frac{8\nu\nu'\pi i}{n}} \frac{\vartheta_{11}\left(u + \frac{2\nu}{n}\right)}{\vartheta_{g_1,g_2}\left(u + \frac{2\nu}{n}\right)},}

taken over $\nu$ through a complete system of residues for the (odd)
modulus $n$. $\nu'$ is an arbitrary integer, and $g_1$, $g_2$ is one of
the three even characteristics $(0,0)$, $(0,1)$, $(1,0)$. This sum
is independent of the particular system of residues that has been
taken for $\nu$.

\begingroup
\footnotesize
\textsuperscript{1}) These relations are due to \gesperrt{Abel} (Oeuvres complètes, Bd. I,
S. 523; Bd. II, S. 251 der neuen Ausgabe). The proof of these relations
given above follows \gesperrt{Hermite} (Crelles Journal, Bd. 32,
S. 283). Also to be mentioned is \gesperrt{Sylow}, Christiania Videnskabsselskabs
Forhandlinger 1864 und 1871. \gesperrt{Kronecker}, Berichte der Berliner Akademie,
19. Juli 1875. \gesperrt{Engel}, Berichte der Sächsischen Gesellschaft der Wissenschaften, 31. Juli 1884.
\endgroup
